Die Motivation für diese Arbeit ergibt sich aus der wachsenden Bedeutung barrierefreier Inhalte im öffentlich-rechtlichen Rundfunk sowie aus den Grenzen bestehender automatisierter Ansätze zur Textvereinfachung. 
Öffentliche rechtliche Medien haben den besonderen gesellschaftlichen Auftrag, Informationen für möglichst alle Bevölkerungsgruppen zugänglich zu machen.
Mit der wachsenden Menge an digitalen Medien steigt zugleich der Bedarf an skalierbaren, automatisierten Verfahren zur Erstellung barrierefreier Texte.\\
In den letzten Jahren haben moderne Large Language Models (LLMs) erhebliche Fortschritte in der Textgenerierung und -umformulierung erzielt.
Sie ermöglichen flexible, kontextabhängige Vereinfachungen und übertreffen klassische regelbasierte Systeme im Bezug auf sprachliche Natürlichkeit.
Die Qualität der generierten Texte hängt stark von Prompt-Formulierungen und Modellparametern ab.
Dies führt zu inkonsistenten Ergebnissen, die für den produktiven Einsatz in redaktionellen Prozessen problematisch sind.

Ein automatisiertes Bewertungssystem verspricht hier einen deutlichen Mehrwert, indem objektive Metriken, klar definierte Qualitätskriterien und eine strukturierte Feedbackschleife eingeführt werden um die Textqualität nachhaltig zu verbessern. \\
Ein solches System ermöglicht es, Schwächen in generierten Texten systematisch zu identifizieren und iterativ auszubessern, um somit die händische Nachbearbeitung zu reduzieren.
Gleichzeitig wird die redaktionelle Arbeit entlastet, da weniger manuelle Nachbearbeitung notwendig ist. \\
Diese Arbeit ist somit sowohl technisch als auch redaktionell motiviert.
Einerseits durch die Weiterentwicklung automatisierter Textvereinfachung, andererseits durch die Verbesserung von Effizienz und Verlässlichkeit im Produktionsprozess.

Ziel ist die Konzeption, Implementierung und Evaluation eines automatisierten Systems zur Bewertung und iterativen Verbesserungen von Texten in Leichter Sprache Plus.
Das System soll in die bestehende LS+ Pipeline integriert werden und generierte Texte anhand definierter linguistischer, semantischer und formaler Kriterien analysieren.
Auf Basis der Bewertung soll automatisiert entschieden werden, ob und in welcher Form eine iterative Verbesserung der Texte erforderlich ist.
Die Qualität, Verständlichkeit und inhaltliche Verlässlichkeit der generierten Texte soll systematisch stabilisiert und verbessert werden.